\chapter{Conservation Laws}

\section{What Should Never Disappear}

Chapter 18 answered a question about single instants: which field states are admissible. Chapters 19 and 20 answered a question about motion: how an admissible state evolves, whether by descent or by conservative oscillation. Neither of these, on its own, asks the question this chapter takes up. Admissibility is a property C either grants or withholds at a given moment; it says nothing about whether some quantity, once present, is guaranteed to remain present as an admissible trajectory unfolds. Conservation is that further claim, and it is a claim about trajectories, not about states considered one at a time. A single snapshot of the field conserves nothing, in the sense this chapter means the word; only a committed sequence of updates, an actual history rather than a single instant of it, can be said to conserve anything at all.

Chapter 20 already supplied one example, in passing, without naming it as an instance of a general pattern. A Hamiltonian trajectory conserves H exactly, by construction, and conserves the symplectic structure binding Φ and v together at every point along its path. This chapter generalizes that observation, and asks what else, across the whole apparatus this book has built since Chapter 8, deserves to be called conserved rather than merely constrained.

\section{Symmetry and Conservation}

The relevant tool is Noether's theorem, applied here to the semantic Lagrangian introduced in Chapter 19. Its content, stated informally, is this: every continuous symmetry of L, every transformation of the field's variables that leaves L's value unchanged, corresponds to a quantity that stays exactly constant along any trajectory satisfying the equations of motion Chapters 19 and 20 developed. A symmetry is not a coincidence to be discovered by inspection alone; it is a structural feature of L, and once identified, it hands the theory a conservation law for free, without needing to verify the conservation independently by tracking the quantity through every possible trajectory.

\section{Three Conserved Quantities}

Three symmetries of this book's framework are worth making explicit, because each corresponds to a conservation law that gives earlier chapters' informal claims a precise form.

The first is time-translation symmetry: nothing about L depends on which update index a trajectory happens to be labeled as starting from, only on the field's actual configuration and its rate of change. This is the symmetry Chapter 20 already exploited without naming it, and its conserved quantity is H itself, energy in this book's sense of the word, exactly constant along any trajectory obeying Hamilton's equations. Chapter 16's identity, refined in Chapter 20 into a persistent symplectic orbit, is now recognizable as this conservation law's signature: an object persists exactly because its trajectory conserves the quantity time-translation symmetry guarantees, and ceases to persist, in this book's sense, exactly when some external event forces it off that conserved orbit entirely.

The second is a domain-relabeling symmetry, closely related to Chapter 15's account of commutative updates. Two commitments to disjoint regions of the field, committed in either order, leave the field's total causal content unchanged, and this reordering is a genuine symmetry of the write cycle: nothing about which commutative update happened first is visible in the resulting history. Its conserved quantity is the total accumulated residue across the field, in the specific sense that no committed event's contribution to R can be created or destroyed by choosing a different, equally valid ordering of commutative commits, only relabeled in when it is said to have occurred. Residue can still grow, every genuine commit adds to it, but it cannot grow or shrink merely as an artifact of reordering, and this is the conservation law underwriting Chapter 15's earlier claim that commutative updates may commit in any order without affecting the final result.

The third is a translation symmetry across the field's domain itself, closest to the mass-consistency invariant this book's supporting technical material lists among its hard constraints. Total coherence, integrated across the whole domain, is conserved except at explicitly licensed sources and sinks: invention, Chapter 9's legitimate act of resolving genuine uncertainty into settled appearance, is a source, adding coherence where entropy previously left none; refusal, Chapter 15's deliberate rejection of a proposal, and unrecovered decay are sinks, removing it. Between these events, coherence does not simply appear or vanish at a point; it moves, in the sense Chapter 7's vector field v was always meant to describe, from where it is to where flow carries it, and the conservation law states precisely that any local increase in Φ must be traceable either to a licensed source or to inflow from a neighboring region, never to nothing at all.

\section{A Worked Example: Three Ways a History Could Lie}

It is worth grounding section 22.3's three laws in a single scenario, because stated abstractly they can look like one diagnostic wearing three names rather than three genuinely independent checks. Consider a courtyard's committed history across three consecutive updates, the fountain, the bench, the wall's crack, and consider three different ways a corrupted or fabricated version of that history could fail, each caught by a different one of the three laws and by none of the others.

Suppose the fountain's water is recorded as moving measurably faster in the third commit than the first, with no proposed forcing term, no wind, no mechanical input, anywhere in the intervening proposals to account for the increase. This is a violation of the first law: H is not constant along the trajectory, and nothing licensed the change. Chapter 21's spatial check would not catch this at all, since each individual commit, examined on its own, is a perfectly well-glued, locally admissible snapshot; the failure is visible only in the comparison across commits, exactly where a conservation law and not a gluing condition is the relevant tool.

Suppose instead the fountain's water is unremarkable, but the courtyard's commit log records two spatially disjoint, genuinely commutative events, a bench moved on one side, a crack inspected on the other, and the total residue accumulated depends on which order the log claims they occurred in, even though Chapter 15 already established that commutative events must be order-independent in their effect on the field. This is a violation of the second law, and it points to a different kind of corruption than the first: not a physically implausible jump in energy, but a commit log that is not actually a faithful causal-DAG record, residue being relabeled rather than merely reordered, which the first law's energy bookkeeping would never notice.

Suppose finally that both of these check out, energy is constant, residue is order-independent, and yet the courtyard's total coherence, Φ, is recorded as higher after the third commit than before it, with no invention event in the residue log and no inflow recorded in the vector field v from any neighboring region. This is a violation of the third law alone, coherence appearing from nowhere, and it is the one most directly resembling old-fashioned hallucination: certainty materializing without an earning event anywhere in the history to account for it.

Each of these three corruptions could occur entirely independently of the other two, and any one of them could occur in a history that passes every one of Chapter 21's spatial checks at every individual instant, since spatial gluing says nothing about what should stay constant, or change only through accounted means, across the sequence. This is the concrete cash value of insisting on three separate conservation laws rather than one: a history that lies about energy, a history that lies about its own causal ordering, and a history that lies about where coherence came from are three different lies, and this framework needs all three checks, run together, to catch all three.

\section{Conservation Along Orbits, Not at Points}

It is worth restating plainly what these three conservation laws are and are not claims about, because the temptation to read them as constraints on individual states, in the manner of Chapter 18's C, would misplace them entirely. C asks whether a given \(M_t\), considered alone, is admissible. These conservation laws ask whether a given trajectory, a committed sequence \(M_t\), \(M_{t+1}\), \(M_{t+2}\), and so on, holds H, total residue, and total licensed coherence constant, or changing only through accounted sources and sinks, across the whole sequence. A single \(M_t\) cannot violate a conservation law by itself, in the same way a single frame of a film cannot violate the law of motion governing the film's characters; violation is only visible across a sequence, exactly as Chapter 20 already implied when it redescribed identity as a property of an orbit rather than a point. This chapter's three conservation laws are, in this sense, elaborations of that same redescription, applied to residue and coherence rather than only to the symplectic structure Chapter 20 introduced first.

\section{What Conservation Buys the Theory}

This gives the framework a second, genuinely different kind of diagnostic from the one Chapter 21 supplied. Chapter 21's cohomological obstruction catches a spatial failure: local patches, checked at a single instant, that cannot be glued into one coherent whole. Conservation catches a temporal failure: a sequence of states, each individually admissible and each individually well-glued in Chapter 21's sense, whose transitions nonetheless violate what should have stayed constant along the way, residue appearing with no committed event to trace it to, coherence increasing at a point with no source or inflow to license it. A hallucinated history can, in principle, pass every one of Chapter 21's spatial checks at every individual instant, and still be caught by this chapter's temporal ones, because gluing correctly at each moment says nothing about whether the moments, taken together as a trajectory, actually obey the conservation laws that any genuine history is bound by. Spatial coherence and temporal conservation are, in this framework, two independent tests, and a fully trustworthy history has to pass both.

\section{Looking Ahead}

This chapter has asked what a single world's trajectory is obligated to conserve as it evolves. It has not asked anything about how one world relates to another, whether two different persistent fields, built from different domains or different histories entirely, can be meaningfully compared, translated into one another, or recognized as instances of the same underlying kind of structure. That is a question about relationships between worlds rather than properties of any one world's trajectory, and it is where Chapter 23 turns next.
